\documentclass[10pt,a4paper]{amsart}
\usepackage[margin=19mm]{geometry}
\usepackage{amsmath,amssymb,mathtools}
\usepackage[scr=rsfs,cal=euler]{mathalfa}
\usepackage{xfrac}
\usepackage[unicode,hidelinks,bookmarks=false]{hyperref}
\newcommand{\ie}{i.e.\ }
\newcommand{\cf}{cf.\ }
\newcommand{\resp}{respectively\ }
\newcommand{\Subsection}{\subsection}
\providecommand{\nobreakdash}{\nobreak\hbox{-}}
\setlength{\emergencystretch}{2em}
\multlinegap0pt
\usepackage{bm}
\usepackage{bbm}
\usepackage{dsfont}
\usepackage{xspace}
\usepackage{cancel}
\usepackage{mathtools}
\usepackage{amsthm}
\makeatletter
\@ifundefined{lemma}{\newtheorem{lemma}{Lemma}}{}
\makeatother
\newcommand{\up}[1]{^{\mathrm{#1}}}
\title[How long is long enough? Finite-horizon approximation of ene]{How long is long enough? Finite-horizon approximation of energy storage scheduling problems}
\author{El\'ea Prat, Richard M. Lusby, Juan Miguel Morales, Salvador Pineda, Pierre Pinson}
\date{Source: arXiv:2411.17463v3 (CC BY 4.0)}
\begin{document}
\maketitle

\noindent\emph{Source and licence.} How long is long enough? Finite-horizon approximation of energy storage scheduling problems. Version arXiv:2411.17463v3 (CC BY 4.0), distributed under the Creative Commons Attribution 4.0 International licence, as recorded in the arXiv licence metadata for this exact version and shown on the abstract page https://arxiv.org/abs/2411.17463v3. Full rights notice: licenses/arxiv-2411.17463-CC-BY-4.0.txt.\par\smallskip

\noindent\textit{Editorial scope.} Counted unit: the Lemma whose source label is \texttt{lemma1} in the article (source0.tex lines 1065--1071, source label \texttt{lemma1}), delivered together with the article's own complete original proof of it (source0.tex lines 1073--1384, one \texttt{proof} environment of 312 lines). No mathematics is written, repaired or re-derived: statement and proof are the article's own text line for line. Required context retained uncounted: eq:s\_calc (lines 1059--1061). Every definition, display and label of those lines is retained unchanged. Omitted from this package and not redistributed: 1--1058, 1062--1064, 1072--1072, 1385--1584. Nothing inside a retained excerpt is omitted or altered: every retained body is delivered exactly as the article's lines are written, so no omission receipt of any kind accompanies this package. Citations: the retained excerpts carry no citation command at all, so no bibliography entry of the article is transcribed and no reference is redistributed. Reference adaptation: none was needed and none was made; every reference inside the delivered unit resolves inside it, so no unresolved reference or citation is produced. Printed slips kept literal: no printed word, symbol, hypothesis or number of the counted statement or of its proof was altered. Layout and class: the article's own class is \texttt{elsarticle} with packages including setspace, hyperref, graphicx, amsmath, amssymb, amsfonts, appendix, algorithm, algpseudocode, subcaption, pgfplots, tikz, tikzscale, makecell; the wrapper is the lane's pinned \texttt{amsart} editorial wrapper at 10pt on A4, which restates the theorem-like environments the retained text needs and adds the article's own macro definitions that the retained text needs (\texttt{\textbackslash up}), with extra wrapper packages (bm, bbm, dsfont, xspace, cancel, mathtools). The delivered numbering is the unit's own. Assets: no retained span contains a figure environment, a graphics-inclusion command, a PDF-page inclusion, a table or a TikZ picture; no image file is copied into this package, the package declares no asset dependency and no image file is redistributed with it. Licence: version arXiv:2411.17463v3 of this article is distributed under the Creative Commons Attribution 4.0 International licence, as recorded in the arXiv licence metadata for this exact version and shown on the abstract page https://arxiv.org/abs/2411.17463v3. A full rights notice is delivered at licenses/arxiv-2411.17463-CC-BY-4.0.txt. Characters: every retained excerpt is pure ASCII, so the delivered engine, XeLaTeX with its default Latin Modern text font, reports no missing character.
\begin{equation} \label{eq:s_calc}
    s_{t_2} = \rho^{t_2 - t_1} s_{t_1} + \Delta t\sum_{t=t_1+1}^{t_2} \rho^{t_2 - t} \left( \eta\up{C} p_t\up{C} - \frac{1}{\eta\up{D}} p_t\up{D} \right), \quad \forall (t_1,t_2) \in \mathcal{T}, t_1 \leq t_2, s_0 = S\up{init}.
\end{equation}

\begin{lemma} \label{lemma1}
    The following is true
    \begin{enumerate}
        \item Given $\mathbf{\overline{x}} \in \overline{\mathcal{X}}$ and $S\up{end}$ , with $\underline{S}_T \leq S\up{end} \leq \overline{S}_T$, there exists $\mathbf{x}^* \in \mathcal{X}^*$ such that ${s}_t^* \leq \overline{s}_t$, $\forall t \in \mathcal{T}$.
        \item Given $\mathbf{\underline{x}} \in \underline{\mathcal{X}}$ and $S\up{end}$ , with $\underline{S}_T \leq S\up{end} \leq \overline{S}_T$, there exists $\mathbf{x}^* \in \mathcal{X}^*$ such that ${s}_t^* \geq \underline{s}_t$, $\forall t \in \mathcal{T}$.
    \end{enumerate}
\end{lemma}

\begin{proof}
    We consider any solution $\mathbf{\overline{x}} \in \overline{\mathcal{X}}$, and any final level $S\up{end}$, with $\underline{S}_T \leq S\up{end} \leq \overline{S}_T$.
    We suppose that for all solutions of $\textbf{F}(\mathcal{T}, \textbf{C}, S\up{end})$, there exists $t \in \mathcal{T}$ such that $s_{t}^* > \overline{s}_t$ and show a contradiction.
    
    Consider any of these solutions, $\mathbf{x}^* \in \mathcal{X}^*$.
    We identify $t_1 \in \mathcal{T}$, such that $s_{t}^* \leq \overline{s}_{t}$, $\forall t < t_1$ and $s_{t_1}^* > \overline{s}_{t_1}$.
    Since $s_T^* = S\up{end}$, $\overline{s}_T = \overline{S}_T$ and $S\up{end} \leq \overline{S}_T$, $\exists \, t_2 \in \mathcal{T}, \, t_2> t_1$ such that $s_{t}^* > \overline{s}_{t}$, $\forall t \in [t_1, t_2-1]$ and $s_{t_2}^* \leq \overline{s}_{t_2}$.

    We build solutions $\mathbf{x}'$ and $\overline{\mathbf{x}}'$, based on $\mathbf{x}^*$ and $\overline{\mathbf{x}}$ respectively, and derive some properties of the problem.

    We first build $\mathbf{x}'$ based on $\mathbf{x}^*$. We modify the solution at $t_1$, decreasing charge if $p\up{C*}_{t_1}>0$ or increasing discharge if $p\up{C*}_{t_1}=0$, which implies $p\up{D*}_{t_1}\geq0$. Therefore, the complementarity constraint at $t_1$ still holds.

    \textit{\textbf{If $p\up{C*}_{t_1}>0$:}} \\ 
    We decrease the charge at $t_1$ by a small quantity ${\epsilon} > 0$, with ${\epsilon} \leq {p}\up{C*}_{t_1}$, such that ${p}\up{C'}_{t_1}={p}\up{C*}_{t_1}-{\epsilon}$ and the bounds on charge are respected.
    We apply \eqref{eq:s_calc} between $t_1$ and $t \in [t_1,t_2-1]$ to obtain the modified state of energy at $t$:
    \begin{equation}
        {s}_t' = \rho^{t - t_1} s_{t_1}' + \Delta t \sum_{k=t_1+1}^{t} \rho^{t - k} \left( \eta\up{C} p_k\up{C'} - \frac{1}{\eta\up{D}} p_k\up{D'} \right) = {s}_t^* - \rho^{t - t_1} \Delta t \, \eta\up{C} {\epsilon}.
    \end{equation}
    We choose ${\epsilon}$ such that the lower bound on the state of energy is respected, i.e. ${\epsilon} \leq \frac{{s}_t^* - \underline{S}}{\rho^{t - t_1} \Delta t \, \eta\up{C}}$, $\forall t \in [t_1,t_2-1]$, which is possible since $\underline{S} \leq \overline{s}_t < s_t^*$, $ \forall t \in [t_1,t_2-1]$. 
    Since ${s}_t^* \leq \overline{S}$, the upper bound is still respected.

    We now modify the solution at $t_2$ to return on the same trajectory, decreasing discharge if $p\up{D*}_{t_2}>0$ or increasing charge if $p\up{D*}_{t_2}=0$ and $p\up{C*}_{t_2}\geq0$. Therefore, the complementarity constraint at $t_2$ still holds.

    \begin{itemize}
        \item \textit{If $p\up{D*}_{t_2}>0$}: We decrease discharge at $t_2$ by $\rho^{t_2 - t_1} \eta\up{C} \eta\up{D} \epsilon$, with ${\epsilon} \leq \frac{{p}\up{D*}_{t_2}}{\rho^{t_2 - t_1} \eta\up{C} \eta\up{D}} $, such that ${p}\up{D'}_{t_2}={p}\up{D*}_{t_2}-\rho^{t_2 - t_1} \eta\up{C} \eta\up{D} \epsilon$ and the bounds on discharge are respected. At $t_2$, the modified state of energy is
        \begin{equation}
            {s}_{t_2}' = \rho {s}_{t_2-1}^* - \rho^{t_2 - t_1} \Delta t \, \eta\up{C} {\epsilon} +\Delta t \, \eta\up{C} {p}\up{C*}_{t_2} - \Delta t \frac{1}{\eta\up{D}} {p}\up{D*}_{t_2} + \Delta t \frac{1}{\eta\up{D}} \rho^{t_2 - t_1} \eta\up{C} \eta\up{D} \epsilon = {s}_{t_2}^*.   
        \end{equation}
        We are back on the same trajectory, which completes the demonstration that the new solution is feasible.
        The change in the objective function compared to $\mathbf{x}^*$ is $\Delta t \, C_{t_1} \epsilon - \Delta t \, \rho^{t_2 - t_1} \eta\up{C} \eta\up{D} C_{t_2} \epsilon$. Since $\mathbf{x}'$ cannot be better than $\mathbf{x}^*$, we derive
        \begin{equation} \label{eq:star_pc_pd}
                C_{t_1} \leq \rho^{t_2 - t_1} \eta\up{C} \eta\up{D} C_{t_2}.
        \end{equation}
        \item \textit{If $p\up{D*}_{t_2}=0$ and $p\up{C*}_{t_2}\geq0$}: We increase charge at $t_2$ by $\rho^{t_2 - t_1} \epsilon$, with ${\epsilon} \leq \frac{\overline{P}\up{C} - p\up{C*}_{t_2}}{\rho^{t_2 - t_1}}$, such that ${p}\up{C'}_{t_2}={p}\up{C*}_{t_2}+\rho^{t_2 - t_1} \epsilon$ and the bounds on charge are respected.
        Note that ${\epsilon} >0$ is guaranteed by $p\up{C*}_{t_2} < \overline{P}\up{C}$. Indeed, since $\overline{s}_{t_2} < {s}_{t_2}^*$, $\overline{s}_{t_2} \geq {s}_{t_2}^*$ would not be possible otherwise.
        At $t_2$, the modified state of energy is
        \begin{equation}
            {s}_{t_2}' = \rho {s}_{t_2-1}^* - \rho^{t_2 - t_1} \Delta t \, \eta\up{C} {\epsilon} +\Delta t \, \eta\up{C} {p}\up{C*}_{t_2} +  \Delta t \, \eta\up{C} \rho^{t_2 - t_1} {\epsilon} - \Delta t \frac{1}{\eta\up{D}} {p}\up{D*}_{t_2} = {s}_{t_2}^*.
        \end{equation}
        We are back on the same trajectory, which completes the demonstration that the new solution is feasible.
        The change in the objective function compared to $\mathbf{x}^*$ is $\Delta t \, C_{t_1} \epsilon - \Delta t \, \rho^{t_2 - t_1} C_{t_2} \epsilon$. Since $\mathbf{x}'$ cannot be better than $\mathbf{x}^*$, we derive
        \begin{equation} \label{eq:star_pc_pc}
            C_{t_1} \leq \rho^{t_2 - t_1} C_{t_2}.
        \end{equation}
    \end{itemize}

   \textit{\textbf{If $p\up{C*}_{t_1}=0$ and $p\up{D*}_{t_1}\geq0$:}} \\ 
   We increase the discharge at $t_1$ by a small quantity ${\epsilon} > 0$, with ${\epsilon} \leq \overline{P}\up{D} - {p}\up{D*}_{t_1}$, such that ${p}\up{D'}_{t_1}={p}\up{D*}_{t_1}+{\epsilon}$ and the bounds on discharge are respected. Note that ${\epsilon} >0$ is guaranteed by $p\up{D*}_{t_1} < \overline{P}\up{D}$. Indeed, since ${s}_{t_1}^* < \overline{s}_{t_1}$, ${s}_{t_1}^* \geq \overline{s}_{t_1}^*$ would not be possible otherwise.
    We apply \eqref{eq:s_calc} between $t_1$ and $t \in [t_1,t_2-1]$ to obtain the modified state of energy at $t$:
    \begin{equation}
        {s}_t' = \rho^{t - t_1} s_{t_1}' + \Delta t \sum_{k=t_1+1}^{t} \rho^{t - k} \left( \eta\up{C} p_k\up{C'} - \frac{1}{\eta\up{D}} p_k\up{D'} \right) = {s}_t^* - \rho^{t - t_1} \Delta t \frac{1}{\eta\up{D}} {\epsilon}.   
    \end{equation}
    We choose ${\epsilon}$ such that the lower bound on the state of energy is respected, i.e. ${\epsilon} \leq  \frac{\eta\up{D}({s}_t^* - \underline{S})}{\rho^{t - t_1} \Delta t}$, $\forall t \in [t_1,t_2-1]$, which is possible since $\underline{S} \leq \overline{s}_t < s_t^*$, $ \forall t \in [t_1,t_2-1]$. 
    Since ${s}_t^* \leq \overline{S}$, the upper bound is still respected.

    We then modify the solution at $t_2$ in a similar way as before.

    \begin{itemize}
        \item \textit{If $p\up{D*}_{t_2}>0$}: We decrease discharge at $t_2$ by $\rho^{t_2 - t_1} \epsilon$, with ${\epsilon} \leq \frac{p\up{D*}_{t_2}}{\rho^{t_2 - t_1}}$, such that ${p}\up{D'}_{t_2}={p}\up{D*}_{t_2}-\rho^{t_2 - t_1}\epsilon$ and the bounds on discharge are respected.
        At $t_2$, the modified state of energy is
        \begin{equation}
            {s}_{t_2}' = \rho {s}_{t_2-1}^* - \rho^{t_2 - t_1} \Delta t \frac{1}{\eta\up{D}} {\epsilon} +\Delta t \, \eta\up{C} {p}\up{C*}_{t_2} - \Delta t \frac{1}{\eta\up{D}} {p}\up{D*}_{t_2} + \Delta t \frac{1}{\eta\up{D}} \rho^{t_2 - t_1} \epsilon = {s}_{t_2}^*.  
        \end{equation}
        We are back on the same trajectory, which completes the demonstration that the new solution is feasible.
        The change in the objective function compared to $\mathbf{x}^*$ is $ \Delta t \, C_{t_1} \epsilon - \Delta t \, \rho^{t_2 - t_1} C_{t_2} \epsilon$. Since $\mathbf{x}'$ cannot be better than $\mathbf{x}^*$, we derive
        \begin{equation} \label{eq:star_pd_pd}
            C_{t_1} \leq \rho^{t_2 - t_1} C_{t_2}.
        \end{equation}
        \item \textit{If $p\up{D*}_{t_2}=0$ and $p\up{C*}_{t_2}\geq 0$}: We increase charge at $t_2$ by $\frac{ \rho^{t_2 - t_1} \epsilon}{\eta\up{C} \eta\up{D}}$, with ${\epsilon} \leq \frac{\eta\up{C} \eta\up{D} (\overline{P}\up{C} - p\up{C*}_{t_2})}{\rho^{t_2 - t_1}}$, such that ${p}\up{C'}_{t_2}={p}\up{C*}_{t_2}+\frac{\rho^{t_2 - t_1} \epsilon}{\eta\up{C} \eta\up{D}}$. The bounds on discharge are respected.
        At $t_2$, the modified state of energy is
        \begin{equation}
            {s}_{t_2}' = \rho {s}_{t_2-1}^* - \rho^{t_2 - t_1} \Delta t \frac{1}{\eta\up{D}} {\epsilon} +\Delta t \, \eta\up{C} {p}\up{C*}_{t_2} +\Delta t \, \eta\up{C} \frac{\rho^{t_2 - t_1} \epsilon}{\eta\up{C} \eta\up{D}} - \Delta t \frac{1}{\eta\up{D}} {p}\up{D*}_{t_2} ={s}_{t_2}^*.
        \end{equation}
        We are back on the same trajectory, which completes the demonstration that the new solution is feasible.
        The change in the objective function compared to $\mathbf{x}^*$ is $ \Delta t \, C_{t_1} \epsilon - \Delta t \frac{\rho^{t_2 - t_1}}{\eta\up{C} \eta\up{D}} C_{t_2} \epsilon $. Since $\mathbf{x}'$ cannot be better than $\mathbf{x}^*$, we derive
        \begin{equation} \label{eq:star_pd_pc}
            \eta\up{C} \eta\up{D} C_{t_1} \leq \rho^{t_2 - t_1} C_{t_2}.
        \end{equation}
    \end{itemize}

    Next, we build $\overline{\mathbf{x}}'$ based on $\overline{\mathbf{x}}$. We modify the solution at $t_1$, decreasing discharge if $\overline{p}\up{D}_{t_1}>0$ or increasing charge if $\overline{p}\up{D}_{t_1}=0$, which implies $\overline{p}\up{C}_{t_1}\geq0$. Therefore, the complementarity constraint at $t_1$ still holds.

    \textit{\textbf{If $\overline{p}\up{D}_{t_1}>0$:}} \\ 
    We decrease the discharge at $t_1$ by a small quantity $\overline{\epsilon} > 0$, with ${\epsilon} \leq \overline{p}\up{D}_{t_1}$, such that $\overline{p}\up{D'}_{t_1} = \overline{p}\up{D}_{t_1}-\overline{\epsilon}$ and the bounds on discharge are respected. 
    We apply \eqref{eq:s_calc} between $t_1$ and $t \in [t_1,t_2-1]$ to obtain the modified state of energy at $t$:
    \begin{equation}
        \overline{s}_t' = \rho^{t - t_1} \overline{s}_{t_1}' + \Delta t \sum_{k=t_1+1}^{t} \rho^{t - k} \left( \eta\up{C} \overline{p}_k\up{C'} - \frac{1}{\eta\up{D}} \overline{p}_k\up{D'} \right)  = \overline{s}_t + \rho^{t - t_1} \Delta t \frac{1}{\eta\up{D}} {\epsilon}.  
    \end{equation}
    We choose $\overline{\epsilon}$ such that the upper bound on the state of energy is respected, i.e. $\overline{\epsilon} \leq  \frac{\eta\up{D} (\overline{S} - \overline{s}_t)}{\rho^{t - t_1} \Delta t}$, $\forall t \in [t_1,t_2-1]$, which we know is possible since $\overline{s}_t < s_t^* \leq \overline{S}$, $ \forall t \in [t_1,t_2-1]$. 
    Since $\overline{s}_t \geq \underline{S}$, the lower bound is still respected.

    We then modify the solution at $t_2$ to return on the same trajectory, decreasing charge if $\overline{p}\up{C}_{t_2}>0$ or increasing discharge if $\overline{p}\up{C}_{t_2}=0$ and $\overline{p}\up{D}_{t_2}\geq0$. Therefore, the complementarity constraint at $t_2$ still holds.

    \begin{itemize}
        \item \textit{If $\overline{p}\up{C}_{t_2}>0$:} We decrease charge at $t_2$ by $\frac{\rho^{t_2 - t_1} \overline{\epsilon}}{\eta\up{C} \eta\up{D}}$, with $\overline{\epsilon} \leq \frac{\eta\up{C} \eta\up{D} \overline{p}\up{C}_{t_2}}{\rho^{t_2 - t_1}}$, such that $\overline{p}\up{C'}_{t_2}=\overline{p}\up{C}_{t_2} - \frac{\rho^{t_2 - t_1} \overline{\epsilon}}{\eta\up{C} \eta\up{D}}$ and the bounds on charge are respected.
        At $t_2$, the modified state of energy is
        \begin{equation}
            \overline{s}_{t_2}' = \rho \overline{s}_{t_2-1} + \rho^{t_2 - t_1} \Delta t \frac{1}{\eta\up{D}} \overline{\epsilon} + \Delta t \, \eta\up{C} \overline{p}\up{C}_{t_2} - \Delta t \, \eta\up{C} \frac{\rho^{t_2 - t_1} \overline{\epsilon}}{\eta\up{C} \eta\up{D}} - \Delta t \frac{1}{\eta\up{D}} \overline{p}\up{D}_{t_2} = \overline{s}_{t_2}.
        \end{equation}
        We are back on the same trajectory, which completes the demonstration that the new solution is feasible. The change in the objective function compared to $\mathbf{x}^*$ is $- \Delta t \, C_{t_1} \overline{\epsilon} + \Delta t \frac{\rho^{t_2 - t_1}}{\eta\up{C} \eta\up{D}} C_{t_2}  \overline{\epsilon}$. Since $\mathbf{x}'$ cannot be better than $\mathbf{x}^*$, we derive
        \begin{equation} \label{eq:overline_pd_pc}
            \eta\up{C} \eta\up{D} C_{t_1} \geq \rho^{t_2 - t_1} C_{t_2}.
        \end{equation}
        \item \textit{If $\overline{p}\up{C}_{t_2}=0$ and $\overline{p}\up{D}_{t_2}\geq 0$:} We increase discharge at $t_2$ by $\rho^{t_2 - t_1} \overline{\epsilon}$, with $\overline{\epsilon} \leq \frac{\overline{P}\up{D} - \overline{p}\up{D}_{t_2}}{\rho^{t_2 - t_1}}$, such that $\overline{p}\up{D'}_{t_2}=\overline{p}\up{D}_{t_2}+\rho^{t_2 - t_1}\overline{\epsilon}$. The bounds on discharge are respected. Note that $\overline{\epsilon} >0$ is guaranteed by $\overline{p}\up{D}_{t_2} < \overline{P}\up{D}$. Indeed, since $\overline{s}_{t_2} < {s}_{t_2}^*$, $\overline{s}_{t_2} \geq {s}_{t_2}^*$ would not be possible otherwise.
        At $t_2$, the state of energy is
        \begin{equation}
            \overline{s}_{t_2}' = \rho \overline{s}_{t_2-1} + \rho^{t_2 - t_1} \Delta t \frac{1}{\eta\up{D}} \overline{\epsilon} + \Delta t \, \eta\up{C} \overline{p}\up{C}_{t_2} - \Delta t \frac{1}{\eta\up{D}} \overline{p}\up{D}_{t_2} -  \Delta t \frac{1}{\eta\up{D}} \rho^{t_2 - t_1}\overline{\epsilon} =\overline{s}_{t_2}.  
        \end{equation}
        We are back on the same trajectory, which completes the demonstration that the new solution is feasible. The change in the objective function compared to $\mathbf{x}^*$ is $- \Delta t \, C_{t_1} \overline{\epsilon} + \Delta t \, \rho^{t_2 - t_1} C_{t_2} \overline{\epsilon}$. Since $\mathbf{x}'$ cannot be better than $\mathbf{x}^*$, we derive
        \begin{equation} \label{eq:overline_pd_pd}
            C_{t_1} \geq \rho^{t_2 - t_1} C_{t_2}.
        \end{equation}
    \end{itemize}

    \textit{\textbf{If $\overline{p}\up{D}_{t_1}=0$ and $\overline{p}\up{C}_{t_1}\geq0$:}} \\ 
    We increase the charge at $t_1$ by a small quantity $\overline{\epsilon} > 0$, with $\overline{\epsilon} \leq \overline{P}\up{C} - \overline{p}\up{C}_{t_1}$, such that $\overline{p}\up{C'}_{t_1}=\overline{p}\up{C}_{t_1}+\overline{\epsilon}$ and the bounds on charge are respected. Note that $\overline{\epsilon} >0$ is guaranteed by $\overline{p}\up{C}_{t_1} < \overline{P}\up{C}$. Indeed, since ${s}_{t_1}^* < \overline{s}_{t_1}$, ${s}_{t_1}^* \geq \overline{s}_{t_1}^*$ would not be possible otherwise.

    We apply \eqref{eq:s_calc} between $t_1$ and $t \in [t_1,t_2-1]$ to obtain the modified state of energy at $t$:
    \begin{equation}
        \overline{s}_t' = \rho^{t - t_1} \overline{s}_{t_1}' + \Delta t \sum_{k=t_1+1}^{t} \rho^{t - k} \left( \eta\up{C} \overline{p}_k\up{C'} - \frac{1}{\eta\up{D}} \overline{p}_k\up{D'} \right)  = \overline{s}_t + \rho^{t - t_1} \Delta t \, {\eta\up{C}} \overline{\epsilon}.
    \end{equation}
    We choose $\overline{\epsilon}$ such that the upper bound on the state of energy is respected, i.e. $\overline{\epsilon} \leq \frac{\overline{S} -\overline{s}_t}{\rho^{t - t_1} \Delta t \, {\eta\up{C}}}$, $\forall t \in [t_1,t_2-1]$, which is possible since $\overline{s}_t < s_t^* \leq \overline{S}$, $ \forall t \in [t_1,t_2-1]$. 
    Since $\overline{s}_t \geq \underline{S}$, the lower bound is still respected.

    We then modify the solution at $t_2$ in a similar way as before.
    \begin{itemize}
        \item \textit{If $\overline{p}\up{C}_{t_2}>0$: } We decrease charge at $t_2$ by $\rho^{t_2 - t_1} \overline{\epsilon}$, with $\overline{\epsilon} \leq \frac{\overline{p}\up{C}_{t_2}}{\rho^{t_2 - t_1}}$, such that $\overline{p}\up{C'}_{t_2}=\overline{p}\up{C}_{t_2}-\rho^{t_2 - t_1}\overline{\epsilon}$ and the bounds on charge are respected.
        At $t_2$, the state of energy is
        \begin{equation}
            \overline{s}_{t_2}' = \rho \overline{s}_{t_2-1} + \rho^{t_2 - t_1} \Delta t \, {\eta\up{C}} \overline{\epsilon} + \Delta t \, \eta\up{C} \overline{p}\up{C}_{t_2} - \Delta t \, \eta\up{C} \rho^{t_2 - t_1}\overline{\epsilon} - \Delta t \frac{1}{\eta\up{D}} \overline{p}\up{D}_{t_2} = \overline{s}_{t_2}.  
        \end{equation}
       We are back on the same trajectory, which completes the demonstration that the new solution is feasible. The change in the objective function compared to $\mathbf{x}^*$ is $- \Delta t \, C_{t_1} \overline{\epsilon} + \Delta t \, \rho^{t_2 - t_1} C_{t_2} \overline{\epsilon}$. Since $\mathbf{x}'$ cannot be better than $\mathbf{x}^*$, we derive
       \begin{equation} \label{eq:overline_pc_pc}
            C_{t_1} \geq \rho^{t_2 - t_1} C_{t_2}.
        \end{equation}
        \item \textit{If $\overline{p}\up{C}_{t_2}=0$ and $\overline{p}\up{D}_{t_2}\geq0$:} We increase discharge at $t_2$ by $\rho^{t_2 - t_1} {\eta\up{C} \eta\up{D}} \overline{\epsilon}$, with $\overline{\epsilon} \leq \frac{\overline{P}\up{D} - \overline{p}\up{D}_{t_2}}{\rho^{t_2 - t_1} \eta\up{C} \eta\up{D}}$, such that $\overline{p}\up{D'}_{t_2}=\overline{p}\up{D}_{t_2}+\rho^{t_2 - t_1} {\eta\up{C} \eta\up{D}}\overline{\epsilon}$. The bounds on discharge are respected.
        At $t_2$, the modified state of energy is 
        \begin{equation}
            \overline{s}_{t_2}' = \rho \overline{s}_{t_2-1} + \rho^{t_2 - t_1} \Delta t \, {\eta\up{C}} \overline{\epsilon} + \Delta t \, \eta\up{C} \overline{p}\up{C}_{t_2} - \Delta t \frac{1}{\eta\up{D}} \overline{p}\up{D}_{t_2} - \Delta t \frac{1}{\eta\up{D}} \rho^{t_2 - t_1} {\eta\up{C} \eta\up{D}} \overline{\epsilon} = \overline{s}_{t_2}.
        \end{equation}
        We are back on the same trajectory, which completes the demonstration that the new solution is feasible. The change in the objective function compared to $\mathbf{x}^*$ is $- \Delta t \, C_{t_1} \overline{\epsilon} + \Delta t \,\rho^{t_2 - t_1} C_{t_2} {\eta\up{C} \eta\up{D}} \overline{\epsilon}$. Since $\mathbf{x}'$ cannot be better than $\mathbf{x}^*$, we derive
        \begin{equation} \label{eq:overline_pc_pd}
            C_{t_1} \geq \rho^{t_2 - t_1} \eta\up{C} \eta\up{D} C_{t_2}.
        \end{equation}
    \end{itemize}
    
    To ensure that $s_{t_1-1}^* \leq \overline{s}_{t_1-1}$, $s_{t_1}^* > \overline{s}_{t_1}$, $s_{t_2-1}^* > \overline{s}_{t_2-1}$, and $s_{t_2}^* \leq \overline{s}_{t_2}$, the possible combinations of charge and discharge at $t_1$ are limited to the following:
    
    \begin{itemize}
        \item At $t_1$, $p_{t_1}\up{C*}>0$ and $\overline{p}_{t_1}\up{D}=0$. At $t_2$, $p_{t_2}\up{D*}=0$ and $\overline{p}_{t_2}\up{C}>0$.
        Then, with \eqref{eq:star_pc_pc} and \eqref{eq:overline_pc_pc}, we get
        \begin{equation} \label{eq:pc_pc_pc_pc}
            C_{t_1} = \rho^{t_2 - t_1} C_{t_2}.
        \end{equation}
        \item At $t_1$, $p_{t_1}\up{C*}>0$ and $\overline{p}_{t_1}\up{D}=0$. At $t_2$, $p_{t_2}\up{D*}>0$ and $\overline{p}_{t_2}\up{C}>0$.
        Then, with \eqref{eq:star_pc_pd} and \eqref{eq:overline_pc_pc}, we get
        \begin{equation} \label{eq:pc_pd_pc_pc}
            C_{t_1} = \rho^{t_2 - t_1} C_{t_2},
        \end{equation}
        and $\eta\up{C} \eta\up{D} = 1$ and/or $C_{t_1} = C_{t_2} = 0$.
        \item At $t_1$, $p_{t_1}\up{C*}>0$ and $\overline{p}_{t_1}\up{D}=0$. At $t_2$, $p_{t_2}\up{D*}>0$ and $\overline{p}_{t_2}\up{C}=0$.
        Then, with \eqref{eq:star_pc_pd} and \eqref{eq:overline_pc_pd}, we get
        \begin{equation} \label{eq:pc_pd_pc_pd}
            C_{t_1} = \rho^{t_2 - t_1} \eta\up{C} \eta\up{D} C_{t_2}.
        \end{equation}
        \item At $t_1$, $p_{t_1}\up{C*}>0$ and $\overline{p}_{t_1}\up{D}>0$. At $t_2$, $p_{t_2}\up{D*}=0$ and $\overline{p}_{t_2}\up{C}>0$.
        Then, with \eqref{eq:star_pc_pc} and \eqref{eq:overline_pd_pc}, we get
        \begin{equation} \label{eq:pc_pc_pd_pc}
            C_{t_1} = \rho^{t_2 - t_1} C_{t_2},
        \end{equation}
        and $\eta\up{C} \eta\up{D} = 1$ and/or $C_{t_1} = C_{t_2} = 0$.
        \item At $t_1$, $p_{t_1}\up{C*}>0$ and $\overline{p}_{t_1}\up{D}>0$. At $t_2$, $p_{t_2}\up{D*}>0$ and $\overline{p}_{t_2}\up{C}>0$.
        Then, with \eqref{eq:star_pc_pd} and \eqref{eq:overline_pd_pc}, we get
        \begin{equation} \label{eq:pc_pd_pd_pc}
            C_{t_1} =  \rho^{t_2 - t_1} C_{t_2},
        \end{equation}
        and $\eta\up{C} \eta\up{D} = 1$ and/or $C_{t_1} = C_{t_2} = 0$.
        \item At $t_1$, $p_{t_1}\up{C*}>0$ and $\overline{p}_{t_1}\up{D}>0$. At $t_2$, $p_{t_2}\up{D*}>0$ and $\overline{p}_{t_2}\up{C}=0$.
        Then, with \eqref{eq:star_pc_pd} and \eqref{eq:overline_pd_pd}, we get
        \begin{equation} \label{eq:pc_pd_pd_pd}
            C_{t_1} = \rho^{t_2 - t_1} C_{t_2},
        \end{equation}
        and $\eta\up{C} \eta\up{D} = 1$ and/or $C_{t_1} = C_{t_2} = 0$.
        \item At $t_1$, $p_{t_1}\up{C*}=0$ and $\overline{p}_{t_1}\up{D}>0$. At $t_2$, $p_{t_2}\up{D*}=0$ and $\overline{p}_{t_2}\up{C}>0$.
        Then, with \eqref{eq:star_pd_pc} and \eqref{eq:overline_pd_pc}, we get
        \begin{equation} \label{eq:pd_pc_pd_pc}
            \eta\up{C} \eta\up{D} C_{t_1} = \rho^{t_2 - t_1} C_{t_2}.
        \end{equation}
        \item At $t_1$, $p_{t_1}\up{C*}=0$ and $\overline{p}_{t_1}\up{D}>0$. At $t_2$, $p_{t_2}\up{D*}>0$ and $\overline{p}_{t_2}\up{C}>0$.
        Then, with \eqref{eq:star_pd_pd} and \eqref{eq:overline_pd_pc}, we get
        \begin{equation} \label{eq:pd_pd_pd_pc}
            C_{t_1} = C_{t_2} = \rho^{t_2 - t_1} C_{t_2},
        \end{equation}
        and $\eta\up{C} \eta\up{D} = 1$ and/or $C_{t_1} = C_{t_2} = 0$.
        \item At $t_1$, $p_{t_1}\up{C*}=0$ and $\overline{p}_{t_1}\up{D}>0$. At $t_2$, $p_{t_2}\up{D*}>0$ and $\overline{p}_{t_2}\up{C}=0$.
        Then, with \eqref{eq:star_pd_pd} and \eqref{eq:overline_pd_pd}, we get
        \begin{equation} \label{eq:pd_pd_pd_pd}
            C_{t_1} = \rho^{t_2 - t_1} C_{t_2}.
        \end{equation}
    \end{itemize}

    Focusing on $\mathbf{x}^*$, we have the following results:
    \begin{enumerate}
        \item If $p_{t_1}\up{C*}>0$ and $p_{t_2}\up{D*}=0$,
        \begin{equation} \label{eq:pc_pc}
            C_{t_1} = \rho^{t_2 - t_1} C_{t_2}.
        \end{equation}
        \item If $p_{t_1}\up{C*}>0$ and $p_{t_2}\up{D*}>0$, 
        \begin{equation} \label{eq:pc_pd}
            C_{t_1} = \rho^{t_2 - t_1} \eta\up{C} \eta\up{D} C_{t_2},
        \end{equation}
        since in \eqref{eq:pc_pd_pc_pc}, \eqref{eq:pc_pd_pd_pc} and \eqref{eq:pc_pd_pd_pd}, either $\eta\up{C} \eta\up{D} = 1$ and/or $C_{t_1} = C_{t_2} = 0$ so this result stands in any case.
        \item If $p_{t_1}\up{C*}=0$ and $p_{t_2}\up{D*}=0$,
        \begin{equation} \label{eq:pd_pc}
           \eta\up{C} \eta\up{D} C_{t_1} =  \rho^{t_2 - t_1} C_{t_2}.
        \end{equation}
        \item If $p_{t_1}\up{C*}=0$ and $p_{t_2}\up{D*}>0$,
        \begin{equation} \label{eq:pd_pd}
           C_{t_1} = \rho^{t_2 - t_1} C_{t_2}.
        \end{equation}
    \end{enumerate}

    Next, we iteratively build a new solution $\mathbf{x}'$ based on $\mathbf{x}^*$, and satisfying point 1 of the lemma, as follows:
    \begin{description}
        \item[$1\up{st}$ step:] Depending on the case, we modify the solution $\mathbf{x}^*$ at $t_1$ and $t_2$ to obtain $\mathbf{x}'$. We explain below how.
        \item[$2\up{nd}$ step:] For $\mathbf{x}'$, we identify $t_1$ and $t_2$, corresponding to the first interval for which the state of energy is strictly greater. If there is no $t_1$, then $s_t' \leq \overline{s}_t$, $\forall t \in \mathcal{T}$ and the procedure terminates.
        \item[$3\up{rd}$ step:] We update $\mathbf{x}^*$ to $\mathbf{x}'$ and go back to the $1\up{st}$ step.
    \end{description}

    We show that this procedure always converges and that the solution built is indeed optimal, which completes our proof.
    We consider separately the 4 cases identified previously.

    \textit{\textbf{1. $p_{t_1}\up{C*}>0$ and $p_{t_2}\up{D*}=0$.}} \\
    We modify the solution at $t_1$ to ${p}\up{C'}_{t_1}={p}\up{C*}_{t_1}-{\epsilon}$ and at $t_2$ to ${p}\up{C'}_{t_2}={p}\up{C*}_{t_2}+\rho^{t_2 - t_1} \epsilon$, with
    \begin{equation} \label{eq:epsilon_1}
        \epsilon = \min \left\{ {p}\up{C*}_{t_1}, \, \left\{\frac{{s}_t^* - \underline{S}}{\rho^{t - t_1} \Delta t \, \eta\up{C}}, \, \forall t \in [t_1,t_2-1]\right\}, \, \frac{\overline{P}\up{C} - p\up{C*}_{t_2}}{\rho^{t_2 - t_1}} \right\},
    \end{equation}
    which we saw previously is a feasible solution, with a change in the objective of
    \begin{equation}
        \Delta t \, C_{t_1} \epsilon - \Delta t \, \rho^{t_2 - t_1} C_{t_2} \epsilon = 0,
    \end{equation}
    with \eqref{eq:pc_pc}. Therefore, this new solution is also optimal.
    
    We show that for this solution, $s_t' > \overline{s}_t$ for fewer time periods than for $\mathbf{x}^*$.
    We evaluate separately the 3 possibilities:
    \begin{itemize}
        \item ${\epsilon} = {p}\up{C*}_{t_1}$. Then $p\up{C'}_t = 0$ and we move to case 3, for which we show below that we can further modify the solution in a way that $s_t' > \overline{s}_t$ for fewer time periods than for $\mathbf{x}^*$.
        \item ${\epsilon} = \min \{ \frac{s_t^* - \underline{S}}{\rho^{t - t_1} \Delta t \, \eta\up{C}}, \, t \in [t_1, t_2-1] \}$. Then for at least one $t \in [t_1, t_2-1]$, $s_t'= \underline{S} \leq \overline{s}_t$, and the result stands.
        \item ${\epsilon} = \frac{\overline{P}\up{C} - p\up{C*}_{t_2}}{\rho^{t_2 - t_1}}$. Then ${p}\up{C'}_{t_2} = \overline{P}\up{C}$. We have $s_{t_2}'- \rho s_{t_2-1}' = s_{t_2}^* - \rho  s_{t_2-1}' = \Delta t \, \eta\up{C} \overline{P}\up{C} \geq \overline{s}_{t_2} - \rho \overline{s}_{t_2-1}$. Since $s_{t_2}^* \leq \overline{s}_{t_2}$, it means that $s_{t_2-1}' \leq \overline{s}_{t_2-1}$ and the result stands.
    \end{itemize}

    \textit{\textbf{2. $p_{t_1}\up{C*}>0$ and $p_{t_2}\up{D*}>0$.}} \\
    We modify the solution at $t_1$ to ${p}\up{C'}_{t_1}={p}\up{C*}_{t_1}-{\epsilon}$ and at $t_2$ to ${p}\up{D'}_{t_2}={p}\up{D*}_{t_2}- \rho^{t_2 - t_1} \eta\up{C} \eta\up{D} \epsilon$, with
    \begin{equation} \label{eq:epsilon_2}
        \epsilon = \min \left\{ {p}\up{C*}_{t_1}, \, \left\{\frac{{s}_t^* - \underline{S}}{\rho^{t - t_1} \Delta t \, \eta\up{C}}, \, \forall t \in [t_1,t_2-1] \right\}, \, \frac{{p}\up{D*}_{t_2}}{\rho^{t_2 - t_1} \eta\up{C} \eta\up{D}} \right\},
    \end{equation}
    which we saw previously is a feasible solution, with a change in the objective of
    \begin{equation}
       \Delta t \, C_{t_1} \epsilon - \Delta t \, \rho^{t_2 - t_1} C_{t_2} \eta\up{C} \eta\up{D} \epsilon = 0,
    \end{equation}
    with \eqref{eq:pc_pd}. Therefore, this new solution is also optimal.
    
    We show that for this solution, $s_t' > \overline{s}_t$ for fewer time periods than for $\mathbf{x}^*$.
    We evaluate separately the 3 possibilities:
    \begin{itemize}
        \item ${\epsilon} = {p}\up{C*}_{t_1}$. Then $p\up{C'}_t = 0$ and we move to case 3, for which we show below that we can further modify the solution in a way that $s_t' > \overline{s}_t$ for fewer time periods than for $\mathbf{x}^*$.
        \item ${\epsilon} = \min \{ \frac{s_t^* - \underline{S}}{\rho^{t - t_1} \Delta t \, \eta\up{C}}, \, t \in [t_1, t_2-1] \}$. Then for at least one $t \in [t_1, t_2-1]$, $s_t'= \underline{S} \leq \overline{s}_t$, and the result stands.
        \item ${\epsilon} = \frac{{p}\up{D*}_{t_2}}{\rho^{t_2 - t_1} \eta\up{C} \eta\up{D}}$. Then ${p}\up{D'}_{t_2} = 0$ and we move to case 1, for which we have shown that we can further modify the solution in a way that $s_t' > \overline{s}_t$ for fewer time periods than for $\mathbf{x}^*$. 
    \end{itemize}

    \textit{\textbf{3. $p_{t_1}\up{C*}=0$ and $p_{t_2}\up{D*}=0$.}} \\
    We modify the solution at $t_1$ to ${p}\up{D'}_{t_1}={p}\up{D*}_{t_1}+{\epsilon}$ and at $t_2$ to ${p}\up{C'}_{t_2}={p}\up{C*}_{t_2}+\frac{\rho^{t_2 - t_1}}{\eta\up{C} \eta\up{D}}\epsilon$, with
        \begin{equation} \label{eq:epsilon_3}
            \epsilon = \min \left\{\overline{P}\up{D} - {p}\up{D*}_{t_1}, \, \left\{  \frac{\eta\up{D}({s}_t^* - \underline{S})}{\rho^{t - t_1} \Delta t}, \, \forall t \in [t_1,t_2-1] \right\} , \, \frac{\eta\up{C} \eta\up{D} (\overline{P}\up{C} - p\up{C*}_{t_2})}{\rho^{t_2 - t_1}} \right\},
        \end{equation}
        which we saw previously is a feasible solution, with a change in the objective of
        \begin{equation}
            \Delta t \, C_{t_1} \epsilon - \Delta t \, C_{t_2} \frac{\rho^{t_2 - t_1}}{\eta\up{C} \eta\up{D}} \epsilon = 0,
        \end{equation}
        with \eqref{eq:pd_pc}. Therefore, this new solution is also optimal.
        
        We show that for this solution, $s_t' > \overline{s}_t$ for fewer time periods than for $\mathbf{x}^*$.
        We evaluate separately the 3 possibilities:
        \begin{itemize}
            \item ${\epsilon} = \overline{P}\up{D} - {p}\up{D*}_{t_1}$. Then $p\up{D'}_{t_1} = \overline{P}\up{D}$. We have $s_{t_1}'- \rho s_{t_1-1}' = s_{t_1}' - \rho s_{t_1-1}^* = - \Delta t \frac{1}{\eta\up{D}} \overline{P}\up{D} \leq \overline{s}_{t_1} - \rho \overline{s}_{t_1-1}$. Since $s_{t_1-1}^* \leq \overline{s}_{t_1-1}$, it means that $s_{t_1}' \leq \overline{s}_{t_1}$ and the result stands.
            \item ${\epsilon} = \min \{ \eta\up{D} \frac{{s}_t^* - \underline{S}}{\rho^{t - t_1} \Delta t}, \, t \in [t_1, t_2-1] \}$. Then for at least one $t \in [t_1, t_2-1]$, $s_t'= \underline{S} \leq \overline{s}_t$, and the result stands.
            \item ${\epsilon} = \frac{\eta\up{C} \eta\up{D} (\overline{P}\up{C} - p\up{C*}_{t_2})}{\rho^{t_2 - t_1}}$. Then ${p}\up{C'}_{t_2} = \overline{P}\up{C}$. We have $s_{t_2}'- \rho s_{t_2-1}' = s_{t_2}^* - \rho s_{t_2-1}' = \Delta t \, \eta\up{C} \overline{P}\up{C} \geq \overline{s}_{t_2} - \rho \overline{s}_{t_2-1}$. Since $s_{t_2}^* \leq \overline{s}_{t_2}$, it means that $s_{t_2-1}' \leq \overline{s}_{t_2-1}$ and the result stands.
        \end{itemize}

    \textit{\textbf{4. $p_{t_1}\up{C*}=0$ and $p_{t_2}\up{D*}>0$.}} \\
    We modify the solution at $t_1$ to ${p}\up{D'}_{t_1}={p}\up{D*}_{t_1}+{\epsilon}$ and at $t_2$ to ${p}\up{D'}_{t_2}={p}\up{D*}_{t_2}-\rho^{t_2 - t_1} \epsilon$, with
        \begin{equation} \label{eq:epsilon_4}
            \epsilon = \min \left\{\overline{P}\up{D} - {p}\up{D*}_{t_1}, \, \left\{ \eta\up{D} \frac{{s}_t^* - \underline{S}}{\rho^{t - t_1} \Delta t}, \, \forall t \in [t_1,t_2-1] \right\} , \, \frac{p\up{D*}_{t_2}}{\rho^{t_2 - t_1}} \right\},
        \end{equation}
        which we saw previously is a feasible solution, with a change in the objective of
        \begin{equation}
            \Delta t \, C_{t_1} \epsilon - \Delta t \, \rho^{t_2 - t_1} C_{t_2} \epsilon = 0,
        \end{equation}
        with \eqref{eq:pd_pd}. Therefore, this new solution is also optimal.
        
        We show that for this solution, $s_t' > \overline{s}_t$ for fewer time periods than for $\mathbf{x}^*$.
        We evaluate separately the 3 possibilities:
        \begin{itemize}
            \item ${\epsilon} = \overline{P}\up{D} - {p}\up{D*}_{t_1}$. Then $p\up{D'}_{t_1} = \overline{P}\up{D}$. We have $s_{t_1}'- \rho s_{t_1-1}' = s_{t_1}' - \rho s_{t_1-1}^* = - \Delta t \frac{1}{\eta\up{D}} \overline{P}\up{D} \leq \overline{s}_{t_1} - \rho \overline{s}_{t_1-1}$. Since $s_{t_1-1}^* \leq \overline{s}_{t_1-1}$, it means that $s_{t_1}' \leq \overline{s}_{t_1}$ and the result stands.
            \item ${\epsilon} = \min \{ \eta\up{D} \frac{{s}_t^* - \underline{S}}{\Delta t}, \, t \in [t_1, t_2-1] \}$. Then for at least one $t \in [t_1, t_2-1]$, $s_t'= \underline{S} \leq \overline{s}_t$, and the result stands.
            \item ${\epsilon} = \frac{p\up{D*}_{t_2}}{\rho^{t_2 - t_1}}$. Then ${p}\up{D'}_{t_2} = 0$ and we move to case 3, for which we have shown that we can further modify the solution in a way that $s_t' > \overline{s}_t$ for fewer time periods than for $\mathbf{x}^*$.
        \end{itemize}

    By induction, starting again the procedure on the new solution, with updated $t_1$ and $t_2$, we can build a solution $\mathbf{x}' \in \mathcal{X}^*$ for which $s_t' \leq \overline{s}_t$, $\forall t \in \mathcal{T}$, which proves point 1 of the lemma.
    
    The proof for point 2 of the lemma is similar.
\end{proof}

\end{document}
